\section{Finite pullback of the original ample line over an arbitrary base ring}

\begin{lemma}[Very ampleness across a closed coefficient immersion]
\label{lem:kn-relative-very-ample-closed-base}
\lean{EnvelopeProperty.Vanishing.RelativeFiniteAmpleClosedBase.isVeryAmple_comp_of_surjective}
\leanok
\uses{lem:kn-projective-embedding-scalar-extension}
Let $\varphi:R\twoheadrightarrow A$ be a surjective homomorphism of
commutative rings, $q:Z\to\operatorname{Spec}A$ a morphism, and $M$ a
line bundle on $Z$. If $M$ is very ample for $q$, then it is very ample
for $\operatorname{Spec}\varphi\circ q$, with the same line bundle.
\end{lemma}
\begin{proof}
Choose a closed immersion $j:Z\hookrightarrow\mathbf P^n_A$ over $A$
and an isomorphism $j^*\mathcal O_{\mathbf P^n_A}(1)\simeq M$.
The native coefficient map $c:\mathbf P^n_A\to\mathbf P^n_R$ is the
base change of $\operatorname{Spec}\varphi$. The latter is a closed
immersion because $\varphi$ is surjective, so $c$ and $jc$ are closed
immersions. Their structure triangle is the original composite
$Z\to\operatorname{Spec}A\to\operatorname{Spec}R$.
The coordinate hyperplane ideal commutes with every coefficient-ring
change. Its dual therefore gives
$c^*\mathcal O_{\mathbf P^n_R}(1)\simeq\mathcal O_{\mathbf P^n_A}(1)$.
Pullback composition and the chosen comparison identify
$(jc)^*\mathcal O_{\mathbf P^n_R}(1)$ with the same $M$, as required.
\end{proof}

\begin{corollary}[Ampleness across a closed coefficient immersion]
\label{cor:kn-relative-ample-closed-base}
\lean{EnvelopeProperty.Vanishing.RelativeFiniteAmpleClosedBase.isAmple_comp_of_surjective}
\leanok
\uses{lem:kn-relative-very-ample-closed-base}
Under the same hypotheses, an ample line for $q$ is ample for the
composite with $\operatorname{Spec}\varphi$.
\end{corollary}
\begin{proof}
The line stays invertible. Choose a positive tensor power which is
very ample for $q$, and apply the preceding lemma to that same power.
This proves the positive-power criterion over $R$ without changing
the line or the exponent.
\end{proof}

\begin{lemma}[The literal coefficient-image factorization]
\label{lem:kn-relative-coefficient-image-factor}
\lean{EnvelopeProperty.Vanishing.RelativeFiniteAmplePullback.coefficientImageStructure_comp}
\leanok
Let $p:X\to\operatorname{Spec}R$ be a morphism and let
$A$ be the image of the induced homomorphism
$R\to\Gamma(X,\mathcal O_X)$. The canonical map
$X\to\operatorname{Spec}\Gamma(X,\mathcal O_X)$ followed by
$\operatorname{Spec}\Gamma(X,\mathcal O_X)\to\operatorname{Spec}A$
defines a morphism $p_A:X\to\operatorname{Spec}A$ satisfying
$\operatorname{Spec}(R\twoheadrightarrow A)\circ p_A=p$.
\end{lemma}
\begin{proof}
The inclusion $A\hookrightarrow\Gamma(X,\mathcal O_X)$ composed with
the surjection $R\twoheadrightarrow A$ is exactly the original
coefficient homomorphism. Contravariant functoriality of the spectrum
and naturality of the canonical global-section morphism therefore give
the asserted equality with $p$.
\end{proof}

\begin{lemma}[The same ample line over the coefficient image]
\label{lem:kn-relative-ample-coefficient-image}
\lean{EnvelopeProperty.Vanishing.RelativeFiniteAmplePullback.isAmple_over_coefficientImage}
\leanok
\uses{lem:kn-relative-coefficient-image-factor,thm:kn-ample-scalar-extension}
If $L$ is ample for $p$, the same line $L$ is ample for $p_A$.
\end{lemma}
\begin{proof}
The square whose top map is the identity of $X$, whose vertical maps
are $p_A$ and $p$, and whose bottom map is
$\operatorname{Spec}(R\twoheadrightarrow A)$ is Cartesian. Indeed,
the bottom map is a monomorphism, and the literal factorization in the
preceding lemma makes this the pullback square of a morphism which
already factors through that monomorphism. Scalar base change preserves
very ampleness, the same positive tensor power, and the actual
invertible line. Its resulting line on $X$ is the pullback along the
identity and hence canonically the original $L$.
\end{proof}

\begin{theorem}[Finite pullback over the original arbitrary coefficient ring]
\label{thm:kn-relative-finite-ample-original-base}
\lean{EnvelopeProperty.Vanishing.RelativeFiniteAmplePullback.isAmple_pullback}
\leanok
\uses{lem:kn-relative-ample-coefficient-image,cor:kn-relative-ample-closed-base,prop:kn-nadel-whole-finite-ample-pullback}
Let $R$ be any commutative ring, $X$ an integral scheme,
$p:X\to\operatorname{Spec}R$, and $L$ an ample line for $p$.
For every finite morphism $\pi:Y\to X$, the actual line $\pi^*L$
is ample for the original composite $\pi p$. The scheme $Y$ may
be disconnected or nonintegral.
\end{theorem}
\begin{proof}
Form the coefficient image $A$ and the literal factorization $p_A$
above. Since $X$ is integral, its global-section ring is a domain,
and its subring $A$ is a domain. The same original $L$ is ample for $p_A$. Finite
pullback ampleness over a domain, applied to $\pi$ and $p_A$, gives
ampleness of $\pi^*L$ on the entire original $Y$ over $A$.
Compose its very ample positive-power embedding with the closed
coefficient immersion $\mathbf P^n_A\hookrightarrow\mathbf P^n_R$.
The closed-base ampleness corollary retains the same original
$\pi^*L$ and the same positive exponent, and the factorization equality
identifies the resulting structure map with $\pi p$. This proves
ampleness over the original $R$ with no hypothesis on the components
of $Y$.
\end{proof}
